\section{Обзор существующих решений}
\label{sec:Chapter2} \index{Chapter2}

В данном разделе рассмотрены подходы к улучшению и восстановлению речи, а также
ключевые компоненты, используемые в схеме Miipher-2: самоконтролируемые
энкодеры и нейросетевые вокодеры. Для каждого класса методов оценивается
степень соответствия требованиям, сформулированным в разделе~\ref{sec:Chapter1}.

\subsection{Классические методы улучшения речи}

Первые методы подавления шума работали в частотной области, оперируя
кратковременным преобразованием Фурье (STFT) $Y(k,m)$ зашумлённого сигнала, где
$k$ --- индекс частотной полосы, $m$ --- индекс кадра. \emph{Спектральное
вычитание} оценивает спектр мощности шума $\hat{P}_n(k)$ на участках паузы и
формирует оценку амплитуды чистого сигнала
\begin{equation}
    |\hat{X}(k,m)| = \sqrt{\max\!\big(|Y(k,m)|^2 - \alpha\hat{P}_n(k),\; 0\big)},
    \label{eq:specsub}
\end{equation}
где $\alpha\ge 1$ --- коэффициент переоценки шума. \emph{Винеровская фильтрация}
строит оптимальный по критерию минимума среднеквадратической ошибки фильтр
\begin{equation}
    H(k,m) = \frac{\xi(k,m)}{1 + \xi(k,m)},
    \qquad
    \xi(k,m) = \frac{P_x(k,m)}{P_n(k,m)},
    \label{eq:wiener}
\end{equation}
где $\xi$ --- априорное отношение сигнал/шум. Методы на основе \emph{MMSE-оценки
амплитудного спектра} дополнительно учитывают статистические модели речи и шума.
Подробное изложение этих методов приведено в~\cite{rabiner_shafer}.

Классические методы вычислительно дёшевы и не требуют обучения, однако опираются
на предположение о стационарности шума и его статистической независимости от
речи. При нестационарных помехах, реверберации и низком SNR они вносят
характерные музыкальные артефакты и не способны восстанавливать утраченные
компоненты сигнала. Требованиям к качеству из раздела~\ref{sec:Chapter1} эти
методы не удовлетворяют, но служат базовым ориентиром.

\subsection{Нейросетевые дискриминативные методы}

С развитием глубокого обучения~\cite{goodfellow_ru} доминирующим стал
\emph{дискриминативный} подход, в котором нейросеть напрямую предсказывает либо
маску, либо чистый спектр. Маскирующие методы оценивают идеальную отношительную
маску (IRM) или комплексную маску (cIRM), применяемую к спектрограмме входа.
Модель DEMUCS~\cite{demucs} работает непосредственно во временной области,
используя свёрточную архитектуру типа U-Net с пропускными связями, и обеспечивает
обработку в реальном времени.

Дискриминативные модели хорошо подавляют аддитивный шум, но при сильных
деградациях склонны к чрезмерному сглаживанию и не восстанавливают подавленные
гармоники, что ограничивает достижимое перцептивное качество.

\subsection{Генеративные методы}

Генеративный подход рассматривает восстановление как задачу условной генерации
чистого сигнала. Модель SEGAN~\cite{segan} впервые применила
генеративно-состязательное обучение к улучшению речи во временной области.
Диффузионные модели, например SGMSE~\cite{sgmse}, формулируют очистку как
обратный стохастический процесс и достигают высокого качества, но требуют
десятков итераций вывода, что нарушает требование по быстродействию.

Отдельный класс образуют методы \emph{восстановления через ресинтез}:
VoiceFixer~\cite{voicefixer} предсказывает чистую мел-спектрограмму, по которой
нейросетевой вокодер синтезирует сигнал. Именно эта идея --- очистка в
компактном промежуточном представлении с последующим вокодингом --- лежит в
основе семейства Miipher и развивается в настоящей работе, с тем отличием, что в
качестве промежуточного представления используется не мел-спектрограмма, а
признаки SSL-энкодера.

\subsection{Самоконтролируемые энкодеры речи}

Самоконтролируемое обучение позволяет извлекать представления из больших объёмов
немаркированной речи. Модель \textbf{wav2vec~2.0}~\cite{wav2vec2} обучается
контрастной задаче на квантованных латентных представлениях.
\textbf{HuBERT}~\cite{hubert} использует предсказание скрытых кластерных
единиц (полученных $k$-средними) по маскированным участкам, что приближает
обучение к задаче маскированного языкового моделирования. \textbf{WavLM}~\cite{wavlm}
расширяет HuBERT, добавляя в обучение моделирование наложенной речи и шумов
(denoising during pre-training), благодаря чему его представления особенно
устойчивы к деградациям --- свойство, ключевое для задачи восстановления.
Закрытый \textbf{USM}~\cite{usm}, применяемый в Miipher-2, использует
квантователь BEST-RQ и обучен на 12~млн часов многоязычного аудио.

Важно, что разные слои этих моделей кодируют разную информацию: согласно
бенчмарку SUPERB~\cite{superb}, нижние слои несут просодико-акустические
признаки, средние --- фонетические, верхние --- семантические. Это
обусловливает необходимость экспериментального подбора слоя для задачи
восстановления (раздел~\ref{sec:Chapter3}). Сравнение открытых энкодеров и USM
приведено в табл.~\ref{tab:encoders}.

\begin{table}[H]
    \caption{Сравнение самоконтролируемых энкодеров речи}
    \label{tab:encoders}
    \centering
    \footnotesize
    \begin{tabularx}{\textwidth}{l c c c X}
        \toprule
        Энкодер & Параметры & Слоёв & Данные предобуч. & Особенности \\
        \midrule
        wav2vec~2.0 & 95 млн & 12 & 960 ч & контрастная задача \\
        HuBERT Base & 95 млн & 12 & 960 ч & предсказание кластеров \\
        WavLM Base+ & 95 млн & 12 & 94 тыс. ч & устойчивость к шуму \\
        WavLM Large & 316 млн & 24 & 94 тыс. ч & устойчивость к шуму \\
        USM (закрыт) & 2 млрд & 32 & 12 млн ч & BEST-RQ, 300+ языков \\
        \bottomrule
    \end{tabularx}
\end{table}

\subsection{Нейросетевые вокодеры}

Вокодер синтезирует волновую форму по компактному представлению.
Авторегрессионный WaveNet обеспечивает высокое качество, но крайне медленен.
\textbf{HiFi-GAN}~\cite{hifigan} --- свёрточный GAN-вокодер с многомасштабным и
многопериодным дискриминаторами; он генерирует сигнал за один проход и работает
существенно быстрее реального времени, что сделало его фактическим стандартом.
\textbf{WaveFit}~\cite{wavefit} --- неавторегрессионный вокодер, основанный на
итерациях с фиксированной точкой; он сочетает преимущества диффузионных и
GAN-подходов и применяется в оригинальном Miipher. \textbf{BigVGAN}~\cite{bigvgan}
--- универсальный вокодер с периодической активацией, демонстрирующий хорошую
обобщаемость на новых дикторах. Сопоставление вокодеров дано в
табл.~\ref{tab:vocoders}.

\begin{table}[H]
    \caption{Сравнение нейросетевых вокодеров}
    \label{tab:vocoders}
    \centering
    \footnotesize
    \begin{tabularx}{\textwidth}{l c c X}
        \toprule
        Вокодер & Параметры & Тип & Характеристика \\
        \midrule
        WaveNet & $\sim$24 млн & авторегресс. & высокое качество, медленный \\
        HiFi-GAN (V1) & 13,9 млн & GAN, 1 проход & быстрый, открытый \\
        WaveFit-5 & $\sim$15 млн & итеративный & 5 итераций, открытый \\
        BigVGAN & 112 млн & GAN, 1 проход & универсальный, тяжёлый \\
        \bottomrule
    \end{tabularx}
\end{table}

\subsection{Преимущества обработки в пространстве признаков}

Обработка в пространстве SSL-признаков, в отличие от обработки во временной или
спектральной области, обладает рядом преимуществ для задачи восстановления.
Во-первых, признаки SSL-энкодера инвариантны ко многим деградациям: благодаря
предобучению на разнообразных данных близкие по содержанию, но по-разному
искажённые сигналы отображаются в близкие представления, что упрощает задачу
очистителя. Во-вторых, частота кадров признаков (50~Гц) на три порядка ниже
частоты дискретизации сигнала, поэтому очиститель работает с короткими
последовательностями и обучается быстро. В-третьих, разделение системы на
очиститель и вокодер позволяет переиспользовать предобученные компоненты и
обучать их по отдельности. Платой за эти преимущества является зависимость от
качества энкодера и невозможность напрямую контролировать фазу сигнала ---
последняя полностью восстанавливается вокодером.

\subsection{Методология оценки качества}

Качество систем восстановления речи оценивают \emph{субъективными} и
\emph{объективными} методами. Субъективная оценка (MOS) наиболее достоверна, но
дорога и плохо воспроизводима. Объективные \emph{интрузивные} метрики (PESQ,
STOI, SI-SDR) сравнивают выход с эталоном и хорошо подходят для контролируемого
эксперимента, однако штрафуют ресинтез за расхождение с эталоном даже при
высоком перцептивном качестве. \emph{Неинтрузивные} метрики (DNSMOS) не требуют
эталона и лучше отражают восприятие генеративных систем. Поскольку
восстановление через ресинтез по природе порождает сигнал, не совпадающий с
эталоном пофазно, в настоящей работе используется совокупность всех трёх типов
метрик, что позволяет получить сбалансированную оценку.

\subsection{Модели Miipher и Miipher-2}

Модель \textbf{Miipher}~\cite{miipher1} реализует схему <<очистка признаков +
вокодер>>: из зашумлённого сигнала извлекаются признаки w2v-BERT, очиститель на
основе DF-Conformer (около 100~млн параметров) предсказывает чистые признаки с
дополнительным текстовым кондиционированием через PnG-BERT, после чего вокодер
WaveFit-5 синтезирует сигнал. Именно с её помощью был создан корпус
LibriTTS-R~\cite{librittsr}.

\textbf{Miipher-2}~\cite{miipher2} упрощает и масштабирует эту схему. Вместо
DF-Conformer применяются \emph{параллельные адаптеры} (около 20~млн параметров),
встраиваемые в замороженный энкодер USM; текстовое и дикторское
кондиционирование исключено полностью, что делает модель многоязычной и
независимой от расшифровки. Вокодер WaveFit оптимизирован по памяти. Авторы
сообщают о коэффициенте реального времени $\mathrm{RTF} = 0{,}0078$ на ускорителе
TPU и возможности обработки порядка миллиона часов аудио за несколько суток;
прирост DNSMOS на LibriTTS составил с $2{,}71$ до $2{,}87$.

\subsection{Выводы по обзору}

Схема Miipher-2 в наибольшей степени отвечает требованиям из
раздела~\ref{sec:Chapter1}: она объединяет устойчивость SSL-представлений к
деградациям, высокое качество ресинтеза и однопроходную (а потому быструю)
работу вокодера. Главным препятствием для воспроизведения является закрытый
энкодер USM. Поэтому в настоящей работе исследуется замена USM открытыми
энкодерами WavLM и HuBERT и подбор подходящего открытого вокодера (HiFi-GAN,
WaveFit). Остающиеся открытыми вопросы --- с какого слоя извлекать признаки и
какая комбинация <<энкодер — вокодер>> оптимальна --- решаются в
разделе~\ref{sec:Chapter3}.
